\documentclass[UTF8,fontset=fandol,11pt]{ctexart}
\usepackage[a4paper,margin=24mm]{geometry}
\usepackage{amsmath,amssymb,amsthm,graphicx,tabularx,booktabs,hyperref}
\newtheorem{definition}{定义}[section]
\newtheorem{theorem}[definition]{定理}
\newtheorem{example}[definition]{例题}
\newtheorem{exercise}[definition]{练习}
\hypersetup{unicode=true,colorlinks=true,linkcolor=blue}
\title{实际表格语法探针}
\date{}
\begin{document}
\maketitle
\par\noindent 离线合成演示：不代表真实模型对本专题的教学验收。\par
\section{方法选择}\label{sec:S05}
% UNIT S05.U02
\begin{theorem}
\label{thm:convergence-method-selection}设 \(\{a_n\}\) 是实数数列。判断其收敛性和计算极限时，以下方法在相应条件下有效：
\[
\begin{array}{c|c|c}
\text{方法}&\text{需要验证的条件}&\text{可得到的结论}\\ \hline
\text{$\varepsilon$--$N$ 定义}&\text{对任意 $\varepsilon>0$ 构造统一的正整数 $N$}&a_n\to L\\
\text{四则运算法则}&\text{各相关数列的极限存在，且商的分母极限非零}&\text{直接计算极限}\\
\text{夹逼定理}&\text{最终成立的上下界具有相同极限}&\text{中间数列收敛到该极限}\\
\text{单调有界原理}&\text{数列单调且有相应方向的界}&\text{数列收敛}.
\end{array}
\]
其中，任何一种方法都不能省略数列的定义域条件；特别是，不能因为某个因子有界就直接把它当作收敛数列使用乘积极限法则。
\end{theorem}
\end{document}
